\chapter{دمانس و سندرم‌های فراموشی}
\sourcepage{۱}{20261005\_015834163.jpg}
\section{آلزایمر}
\begin{itemize}
\item \hl{شایع‌ترین علت دمانس} / \hl{زیر ۶۵ سال نادر است}.
\item بالای ۸۵ سال شیوع ۴۰ درصد / در \hl{زنان} شایع‌تر است.
\item برخی موارد فامیلیال است.
\item کرایتریا تشخیصی:
\begin{itemize}
\item اختلال حافظه: اختلال در یادگیری یا به خاطر آوردن.
\item یک یا چند مورد از: آفازی، آپراکسی، آگنوزی.
\item اختلال عملکرد اجرایی: برنامه‌ریزی، سازمان‌دهی، تفکر انتزاعی.
\item افت شناختی.
\item شروع و پیشرفت تدریجی بیماری.
\item علت آن دلیریوم نباشد.
\item عدم وجود سایر علل درگیری \en{CNS} مثل سکته و پارکینسون.
\end{itemize}
\item پاتولوژی: آتروفی منتشر مغز؛ وجود پلاک‌های سنیل (\unclear{\en{Algentophilic}}) و کلاف‌های نوروفیبریلاری.
\item علائم بالینی:
\begin{itemize}
\item \hl{شروع تدریجی} علائم شناختی، رفتاری.
\item فراموشی.
\item اختلالات هیجانی (زودرنجی، پرخاشگری، تهاجم، تغییر خلق).
\item علائم جسمی: احساس گیجی یا سردرد.
\end{itemize}
\item اختلالات شناختی:
\begin{itemize}
\item اختلال حافظه: \hl{زودرس‌ترین نشانهٔ بیماری} / در ابتدا برای اطلاعات جدید (\en{Early}).
\item اختلال تکلم: کم‌حرفی، پارافازی، گم کردن کلمات، صحبت‌های تکراری؛ اختلالات دیسفازیک (غیرروان، ورنیکه).
\item اختلال در شمردن، اختلال در توانایی ذهنی، اختلال فعالیت‌های روزانه.
\end{itemize}
\end{itemize}

\sourcepage{۲}{20261005\_015844674.jpg}
\section*{۱۲ ـ دمانس؛ ادامهٔ آلزایمر}
\begin{itemize}
\item تظاهرات رفتاری و سایکولوژیک:
\begin{itemize}
\item افسردگی، بی‌قراری، بی‌اشتهایی، اضطراب.
\item آلزایمر ناشی از آسیب مغزی پایه با اختلال شخصیت همراه است.
\item رفتارهای هذیانی و سایکوتیک: توهم بینایی و شنوایی.
\item \hl{سندرم‌های سوءتعبیر} در افراد جوان شایع‌تر؛ اختلال \en{identification}.
\item پرخاشگری، تهاجم، بی‌توجهی، رفتارهای نامناسب جنسی.
\item با شدت دمانس رابطهٔ مستقیم دارد.
\end{itemize}
\item تظاهرات حرکتی:
\begin{itemize}
\item تشنج‌های ژنرالیزه؛ میوکلونوس.
\item برادی‌کینزی و ریجیدیتی اکستراپیرامیدال در مراحل نهایی بیماری.
\item \hl{گشاده شدن پاها روی زمین}، قدم‌های کوتاه.
\item رفلکس‌های دورهٔ نوزادی در مراحل پیشرفتهٔ بیماری.
\item بی‌اختیاری ادرار و مدفوع.
\end{itemize}
\item تصویربرداری: \hl{دیلاتاسیون بطن‌های طرفی}.
\item بیوشیمی آلزایمر:
\begin{itemize}
\item \hl{کاهش ۵۰–۹۰ درصدی استیل‌کولین‌ترانسفراز}.
\item پروتئین \en{tau}.
\end{itemize}
\item درمان:
\begin{itemize}
\item اختلال عملکرد شناختی: \hl{آنتاگونیست گلوتامات}؛ \hl{ممانتین}.
\item \hl{مهارکنندهٔ استیل‌کولین‌استراز}: دونپزیل، ریواستیگمین، گالانتامین.
\item اختلال رفتاری:
\begin{itemize}
\item آنتی‌سایکوتیک‌ها: هالوپریدول، ریسپریدون، کوئتیاپین، تیوتیکسین.
\item ضدافسردگی‌ها: سیتالوپرام، سرترالین، فلوکستین، پاروکستین.
\item ضداضطراب‌ها: کاربامازپین.
\end{itemize}
\end{itemize}
\item بیماران \hl{۴–۱۶ سال} (متوسط ۸ سال) پس از شروع علائم می‌میرند.
\end{itemize}

\sourcepage{۳}{20261005\_015851131.jpg}
\section*{۱۲ ـ دمانس}
\section{اختلال شناختی خفیف (\en{MCI})}
\begin{itemize}
\item اختلال حافظهٔ اخیر.
\item سایر عملکردهای شناختی طبیعی / بیمار زندگی روزمره را دنبال می‌کند.
\item شانس آلزایمر سالانه ۱۲–۱۵ درصد است.
\item \en{RF} تبدیل به آلزایمر: \en{HTN}، تغییر سیگنال مادهٔ سفید مغز در \en{MRI}.
\item افزایش \en{Tau pro}، کاهش آمیلوئید \en{$\beta42$} در \en{CSF}، وجود آلل اپسیلون ۴.
\item آتروفی شدید هیپوکامپ در \en{MRI}.
\item معیارهای تشخیص:
\begin{itemize}
\item اختلال شناختی (غالباً حافظه) که توسط افراد دیگر نیز تأیید شود.
\item اختلال حافظه با توجه به سن و تحصیلات.
\item \en{NL} بودن عملکرد شناختی عمومی.
\item \en{NL} بودن فعالیت روزانه.
\item \en{R/O} دمانس.
\end{itemize}
\end{itemize}
\section{دمانس فرونتوتمپورال (\hl{بیماری پیک})}
\begin{itemize}
\item افتراق از آلزایمر: سن شروع پیک \hl{پایین‌تر است} / اختلالات رفتاری از اختلالات شناختی بارزتر هستند؛ تغییرات شخصیتی.
\item بیماری نادر است / می‌تواند فامیلیال باشد؛ \hl{کروموزوم ۱۷}.
\item \hl{آتروفی لوب‌های فرونتال و تمپورال قدامی}.
\item علائم پاتولوژی:
\begin{itemize}
\item انکلوزیون‌های داخل نورون نقره‌دوست (\hl{\en{Pick Bodies}}).
\item سلول‌های پیک.
\item \en{tau pro} در انکلوزیون.
\end{itemize}
\end{itemize}

\sourcepage{۴}{20261005\_015859892.jpg}
\section*{۱۲ ـ دمانس}
\section{بیماری هانتینگتون}
\begin{itemize}
\item نادر، ارثی؛ \hl{شروع تدریجی با کره و دمانس} (\unclear{نورودژنراتیو}).
\item شروع: \hl{۳۰–۵۰ سالگی} / \hl{\en{AD}} / تکرار توالی \en{CAG} به جای توالی کوتاه در کروموزوم ۴.
\item سابقهٔ مرگ زودرس یکی از فامیل مثل پدر.
\item \hl{آتروفی ژنرالیزهٔ مغز}؛ \hl{دژنرسانس نورون‌های بازال} در ناحیهٔ استریاتوم (\hl{هسته‌های کودیت و پوتامن})؛ یافتهٔ پاتوگنومونیک.
\item علائم: دمانس پیش‌رونده (اختلال حافظه، عملکرد ذهن).
\item اختلال حرکتی: کره؛ تشنج؛ سندرم آکینتیک ـ ریجید؛ دیستونی.
\item اختلال حرکات چشم، بلع و تکلم.
\item مبتلایان به کرهٔ هانتینگتون اختلال شناختی ندارند / اختلال در یادآوری دارند.
\item از زمان شروع علائم به‌طور متوسط \hl{۱۵ سال زنده می‌مانند}.
\item درمان: درمان قطعی ندارد؛ \unclear{پروگنوز را بد می‌کند}.
\item کنترل کره: \hl{بلوک‌کنندهٔ دوپامین}؛ هالوپریدول.
\item \hl{\en{SSRI}}، \hl{آنتی‌سایکوتیک نسل ۲}.
\end{itemize}
\section{تفاوت دمانس کورتیکال و ساب‌کورتیکال}
\begin{itemize}
\item \hl{دمانس ساب‌کورتیکال شایع‌تر}: کندی حرکات، اختلال در توجه و انگیزه، بی‌تفاوتی و افسردگی، علائم درگیری مسیرهای اکستراپیرامیدال.
\item \hl{بیماری هانتینگتون و فلج پیشروندهٔ فوق‌هسته‌ای} موجب دمانس ساب‌کورتیکال.
\item یافتهٔ اصلی دمانس کورتیکال: \hl{اختلال شناختی}.
\item آلزایمر و پیک: خیر؛ دمانس کورتیکال.
\end{itemize}

\sourcepage{۵}{20261005\_015904831.jpg}
\section*{۱۲ ـ دمانس}
\section{دمانس \en{Lewy body}}
\begin{itemize}
\item دومین علت شایع دمانس.
\item در \en{$\tfrac{1}{4}$} افراد سالخوردهٔ مبتلا به دمانس.
\item انکلوزیون‌های گرد و ائوزینوفیلیک (\hl{اجسام لوی}) در داخل \hl{سیتوپلاسم} نورون‌های \hl{کورتکس مغز} و \hl{\en{brain stem}}.
\item اجسام لوی دارای دو \en{pro}: \hl{\en{tau}}، \hl{سینوکلئین آلفا}.
\item اسم در پارکینسون به جسم سیاه داده می‌شود.
\item علائم بالینی:
\begin{itemize}
\item کاهش قدرت شناخت \hl{بدون فراموشی واقعی}.
\item \hl{نوسان در قدرت شناخت}.
\item \hl{توهمات بینایی} شکل‌یافته.
\item \hl{علائم پارکینسونیسم}: برادی‌کینزی و ریجیدیتی.
\end{itemize}
\item درمان: \hl{آنتی‌کولین‌استراز}؛ \hl{ممنوعیت آنتی‌سایکوتیک‌ها} به علت عوارض اکستراپیرامیدال.
\end{itemize}
\section{فلج پیشروندهٔ فوق‌هسته‌ای}
\begin{itemize}
\item \hl{فلج حرکات نگاه عمودی چشم‌ها}.
\item فلج کاذب عضلات عصب‌دهی‌شده توسط مدولا (فلج سودوبولبار).
\item دیزآرتری، ریجیدیتی (درگیری عضلات بالا تنه و گردن) $\pm$ ریجیدیتی اکستراپیرامیدال اندام‌ها.
\item سقوط‌های مکرر؛ نقص در یادآوری اطلاعات بینایی.
\item دمانس ساب‌کورتیکال در ۲۰–۶۰ درصد بیماران: کندی فکر، فقدان حافظه، تغییرات شخصیت.
\item پاتولوژی: مرگ نورون‌ها و گلیوز در مناطق ساب‌کورتیکال به‌ویژه پالیدوم، هستهٔ قرمز، هستهٔ ساب‌تالامیک، شبکهٔ مشبک مغز میانی و جسم سیاه.
\end{itemize}

\sourcepage{۶}{20261005\_015912940.jpg}
\section*{۱۲ ـ دمانس}
\section{هیدروسفالی با فشار طبیعی (\en{NPH})}
\begin{itemize}
\item از علل مهم و قابل درمان دمانس.
\item ایدیوپاتیک، ثانویه به ضربه و خونریزی مغزی.
\item تظاهرات بالینی تدریجی و ظرف چند هفته تا چند ماه تکامل می‌یابند:
\begin{enumerate}[label=\en{\alph*.}]
\item دمانس (اختلال حافظه).
\item آتاکسی و اختلال \en{Gait} (اولین و بارزترین علامت).
\item اختلال کنترل اسفنکتر (بی‌اختیاری ادرار).
\end{enumerate}
\item در مراحل پیشرفته: علائم درگیری لوب فرونتال و پیرامیدال.
\item سردرد و ادم پاپی شایع نیست.
\item تصویربرداری: بزرگ شدن بطن‌های جانبی بدون \en{$\uparrow ICP$}.
\item درمان: شانت بطن به پریتوئن (ونتریکولوپریتونئال).
\end{itemize}
\section{بیماری کروتزفلد ـ جاکوب}
\begin{itemize}
\item دمانس پیشرونده با سیر تحت‌حاد طی هفته‌ها و ماه‌ها.
\item اغلب بعد از یک سال موجب مرگ.
\item شایع‌ترین بیماری پریونی که موجب دمانس می‌شود.
\item سنین شایع ۵۰–۷۰ سالگی.
\item دمانس + میوکلونوس + علائم اکستراپیرامیدال.
\item در \en{EEG}: امواج پریودیک شارپ.
\item در \en{CSF} پروتئین \en{14-3-3} یافت می‌شود.
\item در \en{MRI}: در \en{$T_2$}، سیگنال‌های متراکم در گانگلیون‌های بازال.
\end{itemize}

\sourcepage{۷}{20261005\_015918029.jpg}
\section*{۱۲ ـ دمانس}
\section{ارزیابی دمانس}
\begin{itemize}
\item روتین: الکترولیت‌ها، \en{CBC}، \en{vit $B_{12}$}، تست‌های تیروئیدی، سرولوژی سیفلیس، \en{CT} یا \en{MRI}.
\item اختیاری: \en{ESR}، \en{HIV}، سطح سرمی داروها، سرولوژی لایم، ادرار ۲۴ ساعته جهت فلزات سنگین، آنالیز \en{CSF}، \en{CXR}، \en{EKG}، \en{EEG}، \en{PET-Scan} یا \en{SPECT}.
\end{itemize}
\section{دمانس کاذب}
\begin{itemize}
\item افسردگی مهم‌ترین \en{DDX} دمانس است.
\item در دمانس شکایت به شب بدتر می‌شود ولی افسردگی بدتر نمی‌شود!
\item شکایات جسمی و هایپوکندریازیس در افسردگی شایع‌تر است.
\item شروع دمانس تدریجی و موزیانه است.
\end{itemize}
\section{فراموشی سایکولوژیک و اشکال آن}
\begin{itemize}
\item در افسردگی و اضطراب، به دلیل عدم توجه.
\item آمنزی سایکولوژیک: کاهش جزئی و نامتناسب خاطرات اخیر.
\item بیمار حتی نام خود را فراموش می‌کند.
\end{itemize}
